% TOP_PROBLEM: {"id": 5200014, "title": "Open Problems on Billiards and Geometric Optics", "category_id": 11, "category": {"id": 11, "name": "dynamical_systems", "display_name": "Dynamical Systems", "description": "Problems about long-term behavior of deterministic systems, Hamiltonian dynamics, and periodic orbits.", "slug": "dynamical-systems", "order_index": 11, "created_at": "2026-07-31T15:26:25.671Z"}, "status": "partially_solved"}
% FIRST_POSTED: null
% ENTRY_KIND: statement_only
% Imported from: batch03.tex; UnsolvedMath ID 5200014
\documentclass[11pt]{article}
\usepackage[margin=1in]{geometry}
\usepackage{amsmath,amssymb,mathtools}
\usepackage{fontspec,unicode-math}
\setmainfont{Latin Modern Roman}
\setsansfont{Latin Modern Sans}
\setmathfont{Latin Modern Math}
\usepackage{xurl,hyperref}
\hypersetup{colorlinks=true,urlcolor=blue,linkcolor=blue}
\newcommand{\sref}[2]{\hyperlink{source:#1:#2}{[S#2]}}
\newcommand{\cataloguescope}[1]{\paragraph{Recorded mathematical question.}#1}
\begin{document}
% UnsolvedMath ID 5200014; source catalogue code AMR-051-0014
\setcounter{section}{5200013}
\section{Open Problems on Billiards and Geometric Optics}\label{5200014}
{\small\sffamily UnsolvedMath ID 5200014 · Dynamical Systems}
\subsection{Definitions and mathematical statement}

\paragraph{Shared notation.}$\mathbb N=\{1,2,\ldots\}$, and $\mathbb Z,\mathbb Q,\mathbb R,\mathbb C$ denote the integers, rationals, reals and complex numbers. Any other conventions are specified locally.


An oval is a smooth closed strictly convex plane curve. For a direction $u$, let $I_u$ exchange the two endpoints of each parallel chord, with tangency points fixed. Then $F=I_v\circ I_u$ for the chosen ordered pair of directions. For a pencil through $P$, define the analogous chord involution $I_P$, wherever the construction is defined, and set $F=I_Q\circ I_P$. Conjugacy means $h\circ F\circ h^{-1}$ is the stated circle map for a change of coordinate $h$; a M\"obius transformation is a projective fractional-linear map on $\mathbb RP^1$.
\paragraph{Statement recorded in the supplied source.}
For a planar oval $\gamma$, alternately follow chords in two fixed directions to obtain a circle map $F:\gamma\to\gamma$. If $F$ is conjugate to a rotation for every pair of directions, must $\gamma$ be an ellipse? In the projective version, use pencils through points $P,Q$; if for every $P,Q$ with $PQ$ meeting $\gamma$ the resulting $F$ is conjugate to a Möbius transformation, must $\gamma$ be an ellipse? \sref{5200014}{2}
\subsection{Short English statement}
Does rotation conjugacy for every pair of parallel directions characterize ellipses? Does projective conjugacy for every pair of pencils whose joining line meets the oval also characterize ellipses?
\subsection{Sources}
\begin{description}
\item[\hypertarget{source:5200014:1}{S1}] ULAM.AI and the UnsolvedMath Contributors, \emph{UnsolvedMath: A Curated Collection of Open Mathematics Problems}, record 5200014 (AMR-051-0014). CC BY 4.0. \url{https://huggingface.co/datasets/ulamai/UnsolvedMath}.
\item[\hypertarget{source:5200014:2}{S2}] M. Bialy, C. Fierobe, A. Glutsyuk, M. Levi, A. Plakhov and S. Tabachnikov, \emph{Open problems on billiards and geometric optics} (2021), arXiv:2110.10750v2, Serge Tabachnikov, Problem 2, pp. 10--11. \url{https://arxiv.org/pdf/2110.10750v2}.
\end{description}
\label{end:5200014}
\end{document}
